\chapter{Conclusion} \label{sec:Conclusion}

Praxsmith pulled inspiration from many emergent narrative projects, but primarily
Praxis, which laid the scaffolding for the relation types, agent goals, and general
syntax. Design elements from other research listed in \Cref{sec:BackgroundAndLiteratureReview} were
also used. The result was a domain-specific language which, although flawed in
some ways, showed that a powerful system could be created by combining these
components in a more modern setting.

Both study participants and the author were able to achieve satisfactory results
with this system, and many users reported that they felt as though the system was
capable of more than what they were able to explore in the allotted time. However,
survey bias and response analysis leads the author to believe that the development
experience could be greatly improved with tweaks to the language and underlying
system.

\section{Limitations} \label{sec:Limitations}

\subsection{Incomplete Survey} \label{sec:IncompleteSurvey}

Due to timeline constraints, the survey and experiment trials conducted for this
project provided incomplete results. Only 4 people completed the trials all the
way through, and the average time spent on the project was around 1 hour, despite
a target of 3 to 6 hours.

With an involved project such as a development environment\slash game engine, it
is difficult to find active participants willing to put in the time necessary to
become familiar with a new technology. Further research could offer incentive,
for example monetary, to drive longer participation times. Alternatively, a larger
scale game project could be developed alongside the engine, which would allow for
direct feedback from the authors most involved in the process. Larger projects
like Comme il Faut \cite{mccoy_social_2014} and Versu \cite{nelson_prompter_2014} were
also developed this way, with ``Prom Week'' and ``Blood and Laurels'' being
the associated game projects respectively.

\subsection{Survey Bias} \label{sec:SurveyBias}

This survey came with a significant amount of selection bias, as all respondents
to the preliminary survey had high experience with programming languages,
likely due to the areas this project was advertised. An important part of
this project was how easy the system is to use for writers, which ended up not
being determinable from the selection of survey participants. For future research,
more trials should be run specifically with writers who are not familiar
with programming for a more holistic interpretation.

The user response survey was a limiting factor on this research project. As stated
in \Cref{sec:SurveyBias}, the outreach methods of the survey caused it to be skewed
towards those who were already comfortable with programming, which means the writers
who would be using the tool were not evaluated.

\subsection{Personal Bias} \label{sec:PersonalBias}

As multiple of the trial participants were acquainted with one of the authors to varying
degrees, there is personal bias embedded within the results as well. Casual channels
of communciation (Discord messages) were used during the experimentation stage for
questions and feedback about the systems. Although all participants completed the
surveys as intended, this personal bias cannot be overlooked and should be addressed
in future research by conducting a more repeatable experiment.

\section{Future Work} \label{sec:FutureWork}

\subsection{Improved Studies} \label{sec:ImprovedStudies}

To address the survey limitations enumerated in \Cref{sec:Limitations}, more thoroughly
conducted studies could be used. By offering incentive to participants, broadening the
pool of outreach, and allowing for more time to experiment, survey results would cover
a much more accurate representation of the intended audience.

\subsection{Long-Term Project} \label{sec:LongTermProject}

Another avenue to resolve the evaluation limitations is to trial a longer-term project
using Praxsmith to evaluate how it works in a real-world scenario. As noted in \Cref{sec:Limitations},
Comme il Faut \cite{mccoy_social_2014} and Versu \cite{nelson_prompter_2014} had similar
associated projects, allowing for the developers to build off of their own experiences
as users of the system and more accurately evaluate its capabilities and weaknesses.

\subsection{Engine And Interface Improvements} \label{sec:EngineAndInterfaceImprovements}

There are also many potential improvements to the engine and language that would
make it easier to use from the developer standpoint, as outlined in \Cref{sec:Evaluation}.

Perhaps the foremost technical opportunity for future research is expanding upon the
authorability of Praxsmith by adding either:

\begin{itemize}
    \item A higher level language that writes like a play script, something
    authors are likely familiar with. This is what Prompter does for the Versu engine.
    \cite{nelson_prompter_2014} This would likely give much better results for a
    survey of writers working with this project, and it would still give developers
    the flexibility to realize the full potential of the engine.
    \item An authoring user interface that addresses some of the verbosity and
    bloating observed during usage of the Praxsmith language. The authors of the
    Ensemble Engine recognized this limitation and addressed it with a fully featured
    interface, \cite{samuel_ensemble_2015} so a similar extension is worth looking
    into for Praxsmith, potentially replacing the scripting language altogether.
\end{itemize}
